\section*{RESUMO}
\thispagestyle{empty}

A adoção de ferramentas de Infraestrutura como Código (IaC), como o Terraform, transformou o provisionamento de recursos em ambientes de nuvem modernos liderados por práticas DevOps. Contudo, a identificação precisa de falhas e a interpretação de logs extensos de erro continuam sendo desafios complexos que elevam consideravelmente o tempo médio de recuperação (MTTR) das equipes. Diante deste cenário, o presente projeto de pesquisa tem como objetivo geral propor um método estruturado de interpretação de logs, utilizando Inteligência Artificial Generativa, para identificar a causa raiz e sugerir correções seguras em falhas de provisionamento de infraestrutura. A pesquisa metodologicamente classifica-se como aplicada, de caráter exploratório e descritivo. Os procedimentos técnicos abrangerão revisão bibliográfica aprofundada, análise documental de ferramentas consolidadas e a modelagem conceitual de uma arquitetura de solução. Adicionalmente, será estabelecido um ambiente prático de laboratório para induzir falhas controladas em scripts declarativos, cujos logs resultantes serão sistematicamente submetidos a Modelos de Linguagem de Grande Escala (LLMs). A análise dos dados será qualitativa, avaliando a assertividade da identificação do erro, a aderência à documentação oficial e a viabilidade da refatoração sugerida. Espera-se que a injeção de contexto técnico adequado limite o risco de erros por parte da IA, validando a hipótese de que modelos parametrizados otimizam a triagem operacional. Como contribuição principal, o estudo visa entregar um fluxo de trabalho padronizado e métricas de avaliação que auxiliem engenheiros de confiabilidade na integração confiável de IA em suas rotinas de troubleshooting.

\vspace{0.4cm}

\noindent \textbf{Palavras-chave:}
\par
\noindent Terraform; Modelos de Linguagem; Troubleshooting; AIOps; Provisionamento.
